%Este trabalho está licenciado sob a Licença Creative Commons Atribuição-CompartilhaIgual 4.0 Internacional. Para ver uma cópia desta licença, visite https://creativecommons.org/licenses/by-sa/4.0/ ou envie uma carta para Creative Commons, PO Box 1866, Mountain View, CA 94042, USA.

\chapter{Divisibilidade}

 Dados dois números $a, b \in \Z$, com $b \neq 0$ calculamos a divisão de $a$ por $b$ como mostra a figura abaixo:

  \begin{figure}[H]
   \centering
   \fbox{\includegraphics[width=5cm]{./cap_divisib/figs/divisao}}
   \caption{Representação conjuntos numéricos}
  \end{figure}

 logo, dados $a, b \in \Z$ existem $q, r \in \Z$ tais que $a= bq + r$.

 \vskip0.3cm

 \colorbox{azul}{
 \begin{minipage}{0.9\linewidth}
 \begin{center}
  Um número inteiro $a$ é \textbf{divisível}\index{Divisibilidade} por um número inteiro $b$, se a divisão de $a$ por $b$ tem resto $r= 0$.
 \end{center}
 \end{minipage}}

 \vskip0.3cm

 \begin{exem}
 \begin{itemize}
 \item Todo número par é divisível por $2$.
 \item Todo número terminado em $0$ ou $5$ é divisível por $5$.
 \item Todo número divisível por $2$ e $3$ é também divisível por $6$.
 \end{itemize}
 \end{exem}

 \vskip0.3cm

 \colorbox{azul}{
 \begin{minipage}{0.9\linewidth}
 \begin{center}
  Um número $a \in \N$ diferente de $0$ e de $1$ é \textbf{primo}\index{Números!primos} se for divisível apenas por $1$ e por ele mesmo.
 \end{center}
 \end{minipage}}


 \begin{exem}
 Primos: $\{2, 3, 5, 7, 11, 13, \ldots \}$
 \end{exem}

 \begin{teo}[Teorema Fundamental da Aritmética]
\index{Teorema!fundamental da aritmética}
 Todo número $a \in \N$ diferente de $0$ e de $1$ possui uma decomposição única em números primos, em outras palavras, pode ser escrito como produto de números primos.\index{Fatoração}
 \end{teo}

 \begin{defi}
 Um número $a \in \N$ diferente de $0$ e de $1$ cuja decomposição em primos possui números diferentes de $a$ é chamado de \emph{número composto}.\index{Números!compostos} Neste caso, $1$ e $a$ não são os únicos divisores de $a$.\index{Divisor|seealso{MDC}}\index{Múltiplo|seealso{MMC}}
 \end{defi}

 \begin{exem} De números compostos e suas fatorações em números primos:
 \begin{eqnarray}
  25&=& 5 \cdot 5 \\
  12&=& 2 \cdot 2 \cdot 3 \\
  15&=& 3 \cdot 5 \\
  24&=& 2 \cdot 2 \cdot 2 \cdot 3
 \end{eqnarray}
 \end{exem}

 Para ilustrar o algoritmo utilizado na fatoração dos números naturais iremos utilizá-lo para fatorar os números acima.

 \begin{tabular}{c|c}
  Números a ser fatorado & Números primos em ordem crescente \\
  25 & 5 \\
  5  & 5 \\
  1  & = 5.5  \\
 \end{tabular}

 \begin{multicols}{3}
   \begin{tabular}{c|c}
  12 & 2 \\
   6 & 2 \\
   3 & 3 \\
   1 & = 2.2.3 \\
 \end{tabular}

 \begin{tabular}{c|c}
  15 & 3 \\
   5 & 5 \\
   1 & = 3.5 \\
 \end{tabular}

 \begin{tabular}{c|c}
  24 & 2 \\
  12 & 2 \\
   6 & 2 \\
   3 & 3 \\
   1 & 2.2.2.3 \\
 \end{tabular}
 \end{multicols}

 \begin{obs}
 Um número natural sempre é divisível por todos os seus fatores primos e também pelos produtos de seus fatores primos.
 \end{obs}

 \begin{exem}
 Como $12= 2 \cdot 2 \cdot 3$, temos que seus divisores são:\begin{equation}
D(12)= \{1, 2, 3, 4, 6, 12\}.
\end{equation}
 Como $24= 2 \cdot 2 \cdot 2 \cdot 3$, temos que seus divisores são:\begin{equation}
D(24)= \{1, 2, 3, 4, 6, 8, 12, 24\}.
\end{equation}
 \end{exem}

 \section{MDC e MMC}

 \vskip0.3cm
 \colorbox{azul}{
 \begin{minipage}{0.9\linewidth}
 \begin{center}
  \textbf{MDC - máximo divisor comum}
\index{Máximo divisor comum|see{MDC}}\index{MDC}
  Dados dois números $a, b \in \N$, o máximo divisor comum entre eles, é o maior número natural que divide $a$ e $b$. Se MDC$(a, b)= 1$ então $a$ e $b$ são primos entre si.\index{Números!primos!entre si}
 \end{center}
 \end{minipage}}

 \vskip0.3cm

 \colorbox{azul}{
 \begin{minipage}{0.9\linewidth}
 \begin{center}
 \textbf{MMC - mínimo múltiplo comum}
\index{Mínimo múltiplo comum|see{MMC}}\index{MMC}
 Dados dois números $a, b \in \N$, o mínimo múltiplo comum entre eles, é o menor número natural divisível por $a$ e $b$. Se $a$ e $b$ são primos entre si então MMC$(a, b)= ab$.
 \end{center}
 \end{minipage}}


 \begin{exem}
 \begin{eqnarray}
  \text{MDC}(12, 24)&=& 12 \\
  \text{MMC}(12, 24)&=& 24
 \end{eqnarray}
  \end{exem}

 O cálculo do MMC entre dois números pode ser feito rapidamente com auxílio de seguinte algoritmo, que apresentaremos através de exemplos:

  \begin{exem}
  \begin{enumerate}[a)]
  \item MMC(12, 24):

 \begin{tabular}{c|c}
  12, 24 & 2 \\
   6, 12 & 2 \\
   3,  6 & 2 \\
   3,  3 & 3 \\
   1,  1 &
 \end{tabular}
 $MMC(12, 24) = 2.2.2.3= 24$. Neste caso observe que 24= 2.12.

   \item MMC(9, 10):

   \begin{tabular}{c|c}
    9, 10 & 2 \\
    9, 5  & 3 \\
    3, 5  & 3 \\
    1, 5  & 5 \\
    1, 1  & \\
   \end{tabular}
   $MMC(9, 10)= 2.3.3.5= 90= 9.10$. Neste caso como 9 e 10 não possuem nenhum divisor comum, nesta situação dizemos que eles são primos entre si.

   \item MMC(12, 20):

   \begin{tabular}{c|c}
    12, 20 & 2 \\
     6, 10 & 2 \\
     3,  5 & 3 \\
     1,  5 & 5 \\
     1,  1 & \\
   \end{tabular}
  $MMC(12, 20)= 2.2.3.5= 60 < 12.20= 240$. Este é um caso em que calcular o MMC será uma vantagem, pois o MMC é menor que o produto dos dois números.

  \item MMC(7, 15):

  \begin{tabular}{c|c}
   7, 15 & 3 \\
   7,  5 & 5 \\
   7,  1 & 7 \\
   1,  1 &  \\
  \end{tabular}
  $MMC(7, 15)= 3.5.7= 105$. Observe que neste caso o número 7 é primo, e o número 15 não é um múltiplo de 7, sempre que esta situações ocorrer o MMC entre os números será o produto deles.
  \end{enumerate}

 \end{exem}

 \section{Exercícios}
 \begin{exer}
 (FEPESE - 2013) Um ônibus está circulando com 40 passageiros e tem 4 paradas pela frente até o destino final. Sabe-se que na primeira parada entram 7 e saem 9 passageiros do ônibus. Ainda, a cada parada subsequente o número de passageiros que entram é o dobro do número de passageiros que entraram na parada anterior e o mesmo vale para o número de passageiros que saem. Logo, ao chegar ao destino final, o número de passageiros no ônibus é de:
 \begin{enumerate}[a)]
 \item 5
 \item 10
 \item 15
 \item 20
 \item 30
 \end{enumerate}
 \end{exer}

 \begin{exer}
 (FEPESE - 2013) Uma pessoa deve colocar 312 lâmpadas brancas e 845 lâmpadas amarelas em caixas, de maneira que todas as caixas tenham o mesmo número de lâmpadas, o número de lâmpadas amarelas em cada caixa seja igual e não sobrem lâmpadas fora das caixas. Logo, o número de caixas necessárias para fazer tal divisão é:
  \begin{enumerate}[a)]
  \item 8
  \item 9
  \item 10
  \item 11
  \item 13
  \end{enumerate}
 \end{exer}

 \begin{exer}
 (FEPESE - 2016) João trabalha 5 dias e folga 1, enquanto Maria trabalha 3 dias e folga 1. Se João e Maria folgam no mesmo dia, então quantos dias, no mínimo, passarão para que eles folguem no mesmo dia novamente?
  \begin{enumerate}[a)]
  \item 8
  \item 10
  \item 12
  \item 15
  \item 24
  \end{enumerate}
 \end{exer}

 \begin{exer}
 (FEPESE - 2016) Em uma excursão participam 120 homens e 160 mulheres. Em determinado momento é preciso dividir os participantes em grupos formados somente por homens ou somente por mulheres, de maneira que os grupos tenham o mesmo número de integrantes. Neste caso, o número máximo de integrantes em um grupo é:
  \begin{enumerate}[a)]
  \item 10
  \item 15
  \item 20
  \item 30
  \item 40
  \end{enumerate}
 \end{exer}

 \begin{exer}
 (VUNESP - 2017) Uma papelaria precisa organizar seu estoque de cadernos e, para isso, irá utilizar caixas de papelão, colocando em cada uma delas o mesmo número de cadernos. Se forem colocados 30 cadernos em cada caixa, todas as caixas serão utilizadas e 20 cadernos ficarão de fora, mas, se forem colocados 35 cadernos em cada caixa, todos os cadernos serão encaixotados e 2 caixas não serão utilizadas. Se essa papelaria decidir colocar 40 cadernos em cada caixa, todos os cadernos também serão encaixotados, e o número de caixas necessárias será
  \begin{enumerate}[a)]
  \item 12
  \item 14
  \item 16
  \item 18
  \item 20
  \end{enumerate}
 \end{exer}

 \begin{exer}
 (FEPESE - 2013) Em uma fábrica, a entrega do insumo A ocorre a cada 18 dias e do insumo B, a cada 24 dias. Se ambos os insumos são entregues no dia de hoje, em quantos dias ambos serão novamente entregues simultaneamente?
  \begin{enumerate}[a)]
  \item 32
  \item 44
  \item 60
  \item 72
  \item 144
  \end{enumerate}
 \end{exer}

 \begin{exer}
 (VUNESP - 2017) Se, numa divisão, o divisor e o quociente são iguais, e o resto é 10, sendo esse resto o maior possível, então o dividendo é
  \begin{enumerate}[a)]
  \item 131
  \item 121
  \item 120
  \item 110
  \item 101
  \end{enumerate}
 \end{exer}

 \begin{exer}
 (FGV - 2017) Uma corda de 7 metros e 20 centímetros de comprimento foi dividida em três partes iguais. O comprimento de cada parte é:

  \begin{enumerate}[a)]
  \item 2 metros e 40 centímetros;
  \item 2 metros e 50 centímetros;
  \item 2 metros e 60 centímetros;
  \item 2 metros e 70 centímetros;
  \item 2 metros e 80 centímetros.
  \end{enumerate}
 \end{exer}

 \begin{exer}
 (TJ/SC - 2018) Vanda foi ao consultório médico em uma segunda-feira. O médico disse que ela deveria tomar um comprimido de certo remédio todos os dias, durante 180 dias. Vanda começou a tomar o remédio no mesmo dia da consulta e cumpriu exatamente o que disse o médico.
  O primeiro dia que Vanda NÃO precisou tomar o remédio foi:
  \begin{enumerate}
  \item Uma quarta-feira;
  \item Uma quinta-feira;
  \item Uma sexta-feira;
  \item Um sábado;
  \item Um domingo.
 \end{enumerate}
 \end{exer}

  Gabarito: 1 b); 2 e); 3 c); 4 e); 5 b); 6 d); 7 d); 8 a); 9 d).
